\begin{lemma}[The PE color rule is invariant on old slots]
Let \(A^+\) be obtained by appending arbitrary final values to Tashkinov-series
data \(A\) as in \(\noderef{TashkinovSeriesDataAppend}\).  For every old
index \(i\), selected coloring \(\pi\), vertex \(v\), and color \(\gamma\),
\[
 \operatorname{ParallelExtensionColor}(A^+,i^+,\pi,v,\gamma)
 \quad\Longleftrightarrow\quad
 \operatorname{ParallelExtensionColor}(A,i,\pi,v,\gamma).
\]
\end{lemma}
\begin{proof}
The missing-color clause is identical on both sides because it mentions only
\(\pi,v,\gamma\).  We compare the two history alternatives from
\(\noderef{ParallelExtensionColor}\).

First suppose the enlarged history supplies a latest positive PE index
\(h\leq i^+\).  Its bound gives
\(h.val\leq i.val<i.val+1\leq r+1\), so the old index
\(h_0=\langle h.val,h.val<r+1\rangle\) is defined and \(h=h_0^+\).
By \(\noderef{TashkinovSeriesDataAppend}\), the PE kind, chosen vertex, and
delta at \(h\) are exactly those at \(h_0\).  If \(j\leq i\) is any old PE
slot with chosen vertex \(v\), its image \(j^+\leq i^+\) satisfies the
enlarged maximality clause, hence \(j.val\leq h.val=h_0.val\).  Thus \(h_0\)
is the required latest old slot.  If instead the enlarged no-history clause
holds, applying it to every \(h^+\) proves the old no-history clause.

Conversely, an old latest slot \(h\) maps to \(h^+\), with all three relevant
fields unchanged.  Given any enlarged \(j\leq i^+\), the same numerical-bound
argument produces an old \(j_0\) with \(j=j_0^+\); old maximality then gives
\(j.val=j_0.val\leq h.val\).  In the no-history alternative, convert any
purported enlarged supporting PE slot below \(i^+\) to its old index; the old
no-history clause excludes its chosen vertex.  These conversions prove both
directions of the equivalence.
\end{proof}
